\chapter{Data: Preprocessing}
\label{chap:representation}

\section{Curse of Dimensionality}
\label{sec:representation-the-curse-of-dimensionality}
Input data vector spaces grow in size far more than the actual data space.
For example if one drew a map of the world in 2D there woudl be an extremely expressive dataset.
Adding a height dimension to make a 3D map introduces another infinite dimension.
However, this height dimension is \textit{not} fully explored by the data (buildings are only O(10m) tall, planes fly only O(10km) high, \dots) and so there exists a large empty region in our 3D map that goes off to infinity.
This means that all map elements are somewhat close to each other in the height dimension, which obscures many comparisons, like distance metrics.

A simple grid picture makes the scaling really clear, where $N_{\text{points per axis}}^{N_{\text{features}}}$ samples are needed to fill every grid location.

\figpdf[.78\textwidth]{curse_distance}{For one synthetic sample, nearest and farthest distances become less distinguishable as dimension increases.
This illustrates distance concentration rather than asserting a universal law.}

\section{Standardisation}
\label{sec:representation-scaling-standardisation-and-leakage}
\textbf{Standardisation} is the coordinate-wise transform
\[
    x^{\text{standardised}}_j =
    \frac{
    x_j-\mu_j
    }{
    \sigma_{j}
    }.
\]
It equalises scale but does not remove correlations.
It is also known as the \textit{Z-Score}.
Alternatives include:
min--max scaling:
\[
    x^{\text{minmax}}_j = \frac{x - x_{\text{min}}}{x_{\text{max}} - x_{\text{min}}},
\]
robust median/quantile scaling:
\[
    x^{\text{quintile}}_j = \frac{x - \text{median}(x)}{\text{IQR}(x)},
\]
or a physical unit conversion.


\section{Dimensionality Reduction}
\label{sec:representation-dimensionality-reduction-selection-versus-transformation}
\textbf{Feature selection} keeps original variables; \textbf{feature extraction} creates new coordinates.
Selection usually preserves interpretation; extraction can compress correlated structure more efficiently.
L1 regularisation gives a model-based route to sparse selection in settings where individual parameter weights map meaningfully to input features (\cref{sec:stability-regularisation-l1-and-l2}).

\begin{tabularx}{\linewidth}{@{}L{.20\linewidth}YY@{}}
    \toprule
    Method & What it preserves & Main caution \\
    \midrule
    Filter / wrapper selection & Original variables chosen by a criterion & Selection must occur inside CV \\
    PCA & High-variance linear directions & Variance need not predict labels \\
    Autoencoder & Information useful for reconstruction & Reconstruction need not preserve task signal \\
    t-SNE / UMAP & Local neighbourhood structure for visualisation & Global geometry can be distorted \\
    \bottomrule
\end{tabularx}


\paragraph{Distance In High Dimensions is Useless}
In high dimensions distances can become less discriminative. 
As dimensions increase, the dot product between two vectors converges to one. 
At large dimension count, the distance between any two randomly chosen points becomes nearly identical. 
The ratio of the maximum to minimum distance between points approaches 1, making nearest and farthest neighbors almost equidistant.
Solutons:
\begin{itemize}
    \item \textbf{Manhatten Distance:} use the L1 distance metric,
    \item \textbf{Mahalanobis/Ridge Distance:} handle high-dimensional collinearity and small sample sizes better than standard centroid distances,
    \item \textbf{Pre-Reduce Dimensions:} Apply a reduction in advance of any ownstream operation.
\end{itemize}


\section{PCA}
\label{sec:representation-pca-an-eigenvector-view}
For centred data, Principle Componant Analysis (PCA) diagonalises the covariance matrix and keeps the largest-variance eigen-directions:
\begin{equation}
    \underbrace{\Sigma}_{\text{feature covariance}}
    =
    \underbrace{Q}_{\text{variance directions}}
    \underbrace{\operatorname{diag}(\lambda_1,\ldots,\lambda_{N_{\text{features}}})}_{\text{variance scales}}
    Q^\top
    .
\end{equation}
The eigenstructure reminder is in \cref{subsec:la-eigenvectors-and-eigenvalues}.
Keeping the $N_{\text{components}}$ largest-variance directions gives
\[
    z
    =
    Q_{[:N_{\text{components}}]}^\top
    (x-\mu).
\]
Under squared reconstruction error, this is the optimal linear subspace of that rank for the observed data.
Because PCA ignores labels, low-variance predictive signal can be removed.

\section{Whitening: PCA and ZCA}
\label{sec:representation-whitening-pca-and-zca}
Whitening aims for transformed covariance near the identity.
Using the PCA eigensystem, two common transforms are
\begin{align}
    \underbrace{z_{\mathrm{PCAwhite}}}_{\text{whitened in PCA coordinates}}
    &=
    \underbrace{\operatorname{diag}\!\left((\lambda_i+\epsilon)^{-1/2}\right)}_{\text{equalise variance by direction}}
    Q^\top(x-\mu),\\
    \underbrace{z_{\mathrm{ZCA}}}_{\text{whitened near original orientation}}
    &=
    Q\operatorname{diag}\!\left((\lambda_i+\epsilon)^{-1/2}\right)Q^\top(x-\mu).
\end{align}
The stabiliser $\epsilon$ prevents tiny eigenvalues from exploding noise; exact identity covariance is recovered only in the ideal unstabilised full-rank case.
PCA whitening ends in eigenvector coordinates; ZCA rotates back toward the original coordinate orientation.

\figpdf[.96\textwidth]{pca_zca_whitening}{Synthetic correlated data after centring, PCA whitening and ZCA whitening; both decorrelate/rescale, but their output orientations differ.}

\subsection{Whitening versus Normalisation}
\label{sec:representation-whitening-versus-normalisation}
Whitening targets cross-feature linear correlation as well as scale.
Ordinary standardisation changes each feature independently.
BatchNorm, LayerNorm and RMSNorm also normalise selected statistics rather than generally whitening the representation; see \cref{sec:stability-normalisation-control-the-scale-of-intermediate-representations}.

\section{t-SNE, UMAP and latent visualisation}
\label{sec:representation-t-sne-umap-and-latent-visualisation}
t-SNE and UMAP build low-dimensional views that prioritise neighbourhood relationships rather than globally faithful geometry.
UMAP assumes that local neighbour relations reveal a lower-dimensional structure; t-SNE has a neighbourhood-scale hyperparameter commonly called \textbf{perplexity}.
Use either primarily as an exploratory view, not proof that visual clusters are real physical classes.
Check preprocessing, neighbourhood-scale settings, random seeds and external labels before interpreting a plot.

\begin{tabularx}{\linewidth}{@{}L{.21\linewidth}YY@{}}
    \toprule
    Method & Useful assumption & Failure to test \\
    \midrule
    PCA & High variance contains useful linear structure & Signal lies in low-variance directions \\
    Whitening & Linear correlation/scale are nuisance structure & Small-variance directions are mostly noise \\
    t-SNE / UMAP & Local neighbours reveal meaningful organisation & Global distance/cluster appearance is over-interpreted \\
    Autoencoder & Reconstruction preserves useful information & Rare task-critical signal contributes little to average reconstruction \\
    \bottomrule
\end{tabularx}

